\chapter{INTRODUCTION} % Main chapter title
\label{ChapterIntroduction} % For referencing the chapter elsewhere, use \ref{ChapterIntroduction} 



The report shall not have more than five chapters. The title or heading of these chapters are also fixed except for the Chapter named ‘Actual work’. Students have the flexibility to choose the title for this chapter and its sub chapters. 
\\
Introduction chapter for both BTech and MTech. The chapter shall also have the following four mandatory sub-chapters: 


\section{Context and Background Study}
Introduce the topic and provide relevant background information. Explain why the topic is important and why it warrants further research.  Briefly review the current state of knowledge in the field

\section{Need Analysis}
Systematic process of identifying the needs, problems, or gaps between a current state and a desired state. Explain the disparities between the existing conditions and the desired state in a particular project.

\section{Problem Identification}
The identified problem shall be formulated in a systematic manner and provide clarity to decide on the problem statement and objectives

\section{Problem Formulation}
Under this the reason for choosing the particular problem or title for the project shall be explained along with the thought process that was involved in doing so.



